\section{An Empirical Study of Verifier Rejections}
\label{sec:study}

A rejection confronts the developer with a low-level log, and this section studies what that log records about real rejections and where it falls short.
% 被拒时开发者面对的是一份低层 log，本节研究这份 log 记录了真实 rejection 的什么，又在哪里不够用。
We describe how a developer reads a rejection (\S\ref{sec:study:verifier}), assemble \empiricalset, a dataset of real rejections (\S\ref{sec:study:corpus}), characterize the rejections in it (\S\ref{sec:study:bugs}), and measure the log's localization gap (\S\ref{sec:study:gap}).
% 我们讲开发者怎么读一个 rejection（\S\ref{sec:study:verifier}）、构建真实 rejection 数据集 \empiricalset（\S\ref{sec:study:corpus}）、刻画其中的 rejection（\S\ref{sec:study:bugs}），并度量这份 log 的定位 gap（\S\ref{sec:study:gap}）。

\subsection{Reading a Rejection from the Verifier Log}
\label{sec:study:verifier}

Before the kernel runs an eBPF program, the verifier must prove the program safe.
% 在内核运行一个 eBPF 程序前，verifier 必须证明它安全。
The verifier proves safety one instruction at a time, building for each value the proof that it is used safely, such as a packet pointer staying within bounds, and exploring every path.
% verifier 逐条指令地证明安全，为每个值建立"它被安全使用"的 proof，例如某个 packet pointer 始终在界内，并遍历每条路径。
A program that fails a check is rejected, and the verifier returns a log of the abstract state after each instruction, ending in a one-line terminal error at the rejected instruction.
% 一个没通过某检查的程序会被拒，verifier 返回一份记录每条指令后抽象状态的 log，末尾是被拒指令处的一行 terminal error。

The terminal error names the rejection location, the operation where verification stopped, which the verifier reaches only after the program lost the proof at an earlier instruction.
% 这行 terminal error 标出 rejection location，也就是 verification 停下来的那个 operation，而 verifier 是在程序于更早的指令丢失了 proof 之后才到达这里。
To repair the program so it passes the verifier, the developer must learn which proof the verifier required and where the program failed to keep it, because a fix has to restore that proof.
% 要把程序修好、过掉 verifier，开发者必须搞清楚 verifier 需要哪条 proof、以及程序在哪里没能保住它——因为修复就是要把这条 proof 恢复回来。
The error names only the operation that needed the proof, so the developer reconstructs the required proof and its loss by hand from the trace.
% 而这行错只点出需要 proof 的那个 operation，所以这条 required proof 和它的丢失点都要开发者对着 trace 手工重建。

\begin{figure}[t]
  \begin{lstlisting}[style=bpfix,language=C,
      linebackgroundcolor={\ifnum\value{lstnumber}=3 \color{lstgood!40}\fi}]
// correct source; -O0 lowers the field read into a scalar
static int add_one(int x) { return x + 1; }
int prog(struct xdp_md *ctx) { return add_one(ctx->rx_queue_index) & 1; }
  \end{lstlisting}
  \begin{lstlisting}[style=bpflog]
verifier: R2 invalid mem access 'scalar'
  \end{lstlisting}
  \caption{A compiler-lowering rejection: correct source that the verifier reports as \texttt{invalid mem access `scalar'}.}
  \label{fig:packet}
\end{figure}


\subsection{Reproducing Real Rejections in \empiricalset}
\label{sec:study:corpus}

To study rejections that developers actually hit, we assemble \empiricalset, a dataset of real verifier rejections reproduced under one fixed toolchain.
% 为研究开发者真正会遇到的 rejection，我们构建 \empiricalset：一份在固定 toolchain 下复现的真实 verifier rejection 数据集。
The dataset draws on 936 candidate reports from four sources: Stack Overflow questions, GitHub issues, GitHub fix commits, and kernel selftests.
% 数据集取自四个来源的 936 个候选报告：Stack Overflow 提问、GitHub issue、GitHub 修复 commit、内核 selftest。
We rebuild each candidate and load it into kernel 6.15.11 with clang 18 at log level 2, and 235 of the 936 reject under this build and form \empiricalset.
% 我们重建每个候选并加载进 kernel 6.15.11 + clang 18（log level 2），936 个里有 235 个在这套 build 下被拒，构成 \empiricalset。
The rest are dropped because they do not reject under our toolchain, need a specific environment, or lack the source to rebuild.
% 其余被丢弃，原因是在我们的 toolchain 下不被拒、需要特定环境、或缺少可重建的源码。

Each \empiricalset case carries the faulty source and the developer's own fix, both from its report.
% 每个 \empiricalset case 都带出错源码和开发者自己的 fix，二者都来自其报告。
The source and the fix give the ground truth for the rejection and where its repair landed, which the next two subsections use.
% 源码和 fix 给出该 rejection 及其修复落点的 ground truth，后两小节都会用到。

\subsection{Characterizing the Rejections}
\label{sec:study:bugs}

We characterize each \empiricalset rejection by where its repair lands.
% 我们按每个 \empiricalset rejection 的修复落在哪里来刻画它。
In 191 of the 235 rejections the repair changes the program's source.
% 235 个里有 191 个的修复改动了程序源码。
The other 44 reject correct source and are repaired elsewhere: in the compiler (18), the environment (14), or the verifier (12).
% 另外 44 个拒绝的是正确源码，修复发生在别处：编译器（18）、环境（14）或 verifier（12）。
For example, a context-field read that is provably safe is rejected because \texttt{-O0} lowers it into a scalar.
% 例如一个可证安全的 context 字段读取被拒，是因为 \texttt{-O0} 把它 lower 成了 scalar。
The repair is a compiler flag that leaves the source untouched (Figure~\ref{fig:packet}).
% 它的修复是一个不动源码的编译器 flag（图~\ref{fig:packet}）。
The rest of this subsection examines the 191 program bugs.
% 本小节其余部分考察这 191 个程序 bug。

\begin{figure}[t]
  \begin{lstlisting}[style=bpfix,language=C,
      linebackgroundcolor={\ifnum\value{lstnumber}=3 \color{lstgood!40}\fi}]
// correct source; -O0 lowers the field read into a scalar
static int add_one(int x) { return x + 1; }
int prog(struct xdp_md *ctx) { return add_one(ctx->rx_queue_index) & 1; }
  \end{lstlisting}
  \begin{lstlisting}[style=bpflog]
verifier: R2 invalid mem access 'scalar'
  \end{lstlisting}
  \caption{A compiler-lowering rejection: correct source that the verifier reports as \texttt{invalid mem access `scalar'}.}
  \label{fig:packet}
\end{figure}


A program bug fails a safety check, but the failed check does not name the source mistake.
% 一个程序 bug 没通过某个安全检查，但这个失败的检查并不点出源码错误。
We therefore label each bug with its root cause, the source-level error the developer made.
% 因此我们给每个 bug 标注它的 root cause，即开发者犯的源码级错误。
The 191 bugs fall into 12 root causes (Table~\ref{tab:mechanisms}), most specific to eBPF, such as proving a packet bound or pairing a dynptr with its lifetime.
% 这 191 个 bug 落在 12 类 root cause（Table~\ref{tab:mechanisms}），多数是 eBPF 特有的，例如证明一个 packet bound、或把 dynptr 与其生命周期配对。
Of the 12 root causes, only two also occur outside eBPF: missing null checks and out-of-bounds indexing.
% 12 类 root cause 里只有两类在 eBPF 之外也会出现：缺失的 null check 和越界索引。
Together these two cover 34 of the 191 bugs.
% 这两类合计覆盖 191 个 bug 里的 34 个。

\begin{tcolorbox}[colback=blue!5!white,colframe=gray!75!black,left=1mm, right=1mm, top=0.5mm, bottom=0.5mm, arc=1mm]
\textbf{Takeaway \#1:} Most rejections (81\%) are program bugs, but 19\% reject correct source due to compiler, environment, or verifier issues. Of 12 root causes, 10 are eBPF-specific, so repair requires domain knowledge.
% \textbf{要点 \#1：}大多数 rejection（81\%）是程序 bug，但 19\% 拒绝的是正确源码（编译器、环境或 verifier 问题）。12 类 root cause 中 10 类是 eBPF 特有的，因此修复需要领域知识。
\end{tcolorbox}

\begin{table}[t]
  \centering
  \caption{The 191 developer bugs fall into 12 root-cause categories. A category is the specific source-level error the developer made; the \emph{Cases} column gives how many bugs each accounts for.}
  \label{tab:mechanisms}
  \small
  \begin{tabular}{@{}l r@{}}
    \toprule
    Root-cause category & Cases \\
    \midrule
    Unclamped scalar used as offset or length   & 24 \\
    Corrupted or stale dynptr object            & 23 \\
    Packet access without a bound on every path  & 22 \\
    Missing null check                           & 19 \\
    Pointer type or provenance mismatch          & 16 \\
    Unverified address dereferenced              & 16 \\
    Index exceeds object capacity                & 15 \\
    Context or contract misuse                   & 15 \\
    Unpaired resource reference                  & 15 \\
    Interrupt flag not restored in order         & 11 \\
    Probe signature mismatched with the ABI      & 9  \\
    Oversized or uninitialized stack buffer      & 6  \\
    \midrule
    Total                                        & 191 \\
    \bottomrule
  \end{tabular}
\end{table}

\subsection{Measuring the Log's Localization Gap}
\label{sec:study:gap}

With the root causes labeled, we find the terminal error too coarse to name which one a rejection hit.
% root cause 标好之后，我们发现 terminal error 太粗，说不出一个 rejection 撞上的是其中哪一类。
The errno it carries does little to narrow the cause, as \texttt{EINVAL} alone tags 47\% of all rejections.
% 它携带的 errno 几乎无助于缩小原因，单是 \texttt{EINVAL} 就标注了全部 rejection 的 47\%。
The error string is just as coarse, because one string maps to many root causes.
% 这行错的字符串同样粗，因为一个字符串对应许多 root cause。
To measure how often, we mask the registers and offsets in each string and group the 235 rejections by the normalized template.
% 为度量这有多频繁，我们把每个字符串里的寄存器和 offset 掩掉，再按归一化后的 template 给 235 个 rejection 分组。
This collapses 167 distinct strings into 82 templates.
% 这把 167 条不同字符串收成 82 个 template。
Fifteen of the 82 templates each map to more than one root cause.
% 82 个 template 里有 15 个各自对应不止一个 root cause。
The most frequent alone covers nine root causes across 28 rejections (Table~\ref{tab:fanout}).
% 单是最频繁的那一个就跨了 9 类 root cause、28 个 rejection（Table~\ref{tab:fanout}）。

\begin{table}[t]
  \centering
  \caption{The four most frequent message templates, with the cases sharing each template and the distinct root-cause categories they span. A template normalizes registers and offsets to placeholders.}
  \label{tab:fanout}
  \small
  \begin{tabular}{@{}l r r@{}}
    \toprule
    Terminal message template & Cases & Categories \\
    \midrule
    \texttt{R\# invalid mem access `scalar'}   & 28 & 9 \\
    \texttt{invalid access to packet}          & 26 & 5 \\
    \texttt{invalid access to map value}       & 18 & 4 \\
    \texttt{R\# !read\_ok}                     & 13 & 4 \\
    \bottomrule
  \end{tabular}
\end{table}

Because one error matches many root causes, the error alone cannot identify which proof the program lost, or where, the two facts a repair needs.
% 因为一行错对应许多 root cause，单凭这行错无法判断程序丢失的是哪条 proof、又是在哪丢的，而这正是修复需要的两件事。

\begin{tcolorbox}[colback=blue!5!white,colframe=gray!75!black,left=1mm, right=1mm, top=0.5mm, bottom=0.5mm, arc=1mm]
\textbf{Takeaway \#2:} The terminal error is too coarse to guide repair: 47\% of rejections return \texttt{EINVAL}, and one error string maps to as many as nine root causes.
% \textbf{要点 \#2：}terminal error 太粗，无法指导修复：47\% 的 rejection 返回 \texttt{EINVAL}，一个 error 字符串最多对应 9 种 root cause。
\end{tcolorbox}

The trace, however, holds more than the terminal error.
% 但 trace 保留的不止是 terminal error 这一行。
Under \texttt{log\_level=2} the verifier prints its abstract state after every instruction, so the lost proof survives in the printed states, and \S\ref{sec:design} reconstructs the proof from those states to localize the loss.
% 在 \texttt{log\_level=2} 下，verifier 在每条指令后打印它的抽象状态，于是丢失的 proof 留在这些打印状态里，而第 3 节就从这些状态重建该 proof、定位丢失点。